\documentclass{article}
\usepackage{amsmath}
\usepackage{amsthm}
\usepackage{amsfonts}
\usepackage{amssymb}
\usepackage{bbold}
\usepackage{dsfont}
\usepackage{graphicx}
\usepackage{cleveref}
\usepackage{bm}
\usepackage{booktabs}
\usepackage{siunitx}
\usepackage{xcolor}
\usepackage[utf8]{inputenc}
\usepackage[english]{babel}
\newtheorem{definition}{Definition}
\newtheorem{example}{Example}
\newcommand{\exP}[1]{\mathrm{e}^{#1}}
\newcommand{\dif}{\mathrm{d}}
\renewcommand{\vec}[1]{\mathbf{#1}}
\newcommand{\pdagger}{{\phantom{\dagger}}}
\newcommand{\state}[1]{\vert #1\rangle}
\newcommand{\cstate}[1]{\langle #1\vert}
\newcommand{\tr}[1]{\mathrm{tr}\left\{#1\right\}}
\newcommand{\trs}[2]{\, \mathrm{tr}_{#1}\{#2\}}
\newcommand{\expc}[3]{\langle #1\vert #2 \vert #3\rangle}
\usepackage[doi=false,url=false,isbn=false,backend=bibtex,style=nature,natbib=true,sorting=none]{biblatex}
\bibliography{bib.bib}
\usepackage{geometry}
\geometry{
	a4paper,
	total={170mm,257mm},
	left=20mm,
	top=20mm,
}
\begin{document}
\title{Hamiltonian und Nomenklatur in den Dateien}
\author{Rudolf Smorka}
\maketitle
\section{Hamiltonian}
System-Hamiltonian:
\begin{equation}
	\hat{H}_\mathrm{s} = \hat{H}_\mathrm{el} + \hat{H}_\mathrm{el-vib} + \hat{H}_\mathrm{vib}.  \label{eq:H_s}
\end{equation}
mit reinem elektronischen Teil
\begin{equation}
	\begin{split}
		\hat{H}_\mathrm{el} &= \sum_{j,\sigma} \varepsilon_{j\sigma}^\pdagger \hat{d}_{j\sigma}^\dagger \hat{d}_{j\sigma}^\pdagger - t_0 \sum_{j\sigma} (\hat{d}_{j\sigma}^\dagger \hat{d}_{j+1,\sigma}^\pdagger + \mathrm{h.c.}) +\\& \lambda_0^\pdagger\sum_{j\sigma\sigma'} (i\hat{d}_{j\sigma}^\dagger (\vec{v}_j^{(+)}\cdot \bm{\sigma})_{\sigma\sigma'}^\pdagger \hat{d}_{j+2,\sigma'}^\pdagger + \mathrm{h.c.}).\label{eq:H_el}
	\end{split}
\end{equation}
und der vibronischen Teil
\begin{align}
	\hat{H}_\mathrm{el-vib} &= - t_1^\pdagger \sum_{j,\sigma}  \hat{d}_{j\sigma}^\dagger \hat{d}_{j+1\sigma}^\pdagger (\hat{a}_{j}^\dagger + \hat{a}_{j}^\pdagger) \\ & + \lambda_1^\pdagger\sum_{j,\sigma,\sigma'} (i \hat{d}_{j\sigma}^\dagger (\vec{v}_j^{(+)} \cdot \bm{\sigma})_{\sigma\sigma'}^\pdagger \hat{d}_{j+2\sigma'}^\pdagger +\mathrm{h.c.})(\hat{a}_{j}^\dagger + \hat{a}_{j}^\pdagger).\label{eq:H_el-vib}
\end{align}
Hier, $\vec{v}^{(+)}_j = \vec{d}_i\times \vec{d}_{i+1}$, und $\vec{d}_{i+s} = \frac{\vec{r}_i - \vec{r}_{i+s}}{\vert \vec{r}_i - \vec{r}_{i+s}\vert}$ für $\vec{r}_i = (a\cos (\phi_i), a\sin(\phi_i),c_i)$ mit $\phi_i = \frac{2\pi(i-1)M}{\mathbb{M}-1}$ und $c_i = c\frac{i-1}{\mathbb{M}-1}$. Desweiteren sind $a=1$ und $c=1$. 
Der Schwingungshamiltonian ist
\begin{equation}
	\hat{H}_\mathrm{vib} = \sum_{j} \omega_{j}^\pdagger \hat{a}_{j}^\dagger \hat{a}_{j}^\pdagger.
	\label{eq:H_vib} 
\end{equation}
Das elektronische Bad ist
\begin{equation}
	\hat{H}_\mathrm{env} = \sum_{\ell k\sigma} \varepsilon_{\ell k\sigma}^\pdagger c_{\ell k\sigma}^\dagger c_{\ell k\sigma}^\pdagger \label{eq:LeadHamiltonian}
\end{equation}
und der System-Bad Hamiltonian
\begin{align}
	H_\mathrm{s-env} &= \sum_{\sigma\sigma'}(\nu_{L\sigma\sigma'}^\pdagger c_{N_L\sigma}^\dagger d_{1\sigma'}^\pdagger + \nu_{R\sigma\sigma'}^\pdagger c_{1_R\sigma}^\dagger d_{\mathbb{M}\sigma'}^\pdagger + \mathrm{h.c.}). \label{eq:LeadSysHamiltonian}
\end{align} 
Das Bad ist im wide-band-limit zu betrachten. Die Parameter für den obigen Hamiltonian sind:
\begin{align*}
	\varepsilon_0 &= \SI{-240}{\milli\electronvolt}\\
	\omega_0 &= \SI{2}{\milli\electronvolt}\\
	t_0 &= \SI{40}{\milli\electronvolt}\\
	t_1 &= \SI{4}{\milli\electronvolt}\\
	\lambda_0 &= \SI{1}{\milli\electronvolt}\\
	\lambda_1 &= \SI{0.1}{\milli\electronvolt}\\
	\Gamma &= \SI{10}{\milli\electronvolt} (=2\pi \sum_k \vert \nu_k\vert^2\delta(\varepsilon-\varepsilon_k) )\\
	T &= \SI{25}{\milli\electronvolt}\\
	V &= \{0.384,0.44,0.52\}\SI{}{\volt}.
\end{align*}
\section{Hamiltonian Dateien}
\subsection{*EL*}
Die Nomenklatur in den Dateien ist am einfachsten an einem Beispiel zu verstehen. Nehme das System mit drei dots in \texttt{hamiltonian\_1x3\_tPh\_0.0\_thPd\_0.0\_SOC0.001\_SOCPh0.0\_V0.52\_w00.002\_EL\_real.dat}. Der Realteil des rein elektronischen Hamiltonians ist 
\begin{equation}
\mathrm{Re}(H_\mathrm{el}) = \begin{pmatrix}
	& d_{1\uparrow} & d_{1\downarrow} & d_{2\uparrow} & d_{2\downarrow} & d_{3\uparrow} & d_{3\downarrow}\\ \hline
	d_{1\uparrow}^\dagger | & -0.24 &  0.0     &-0.04  & 0.0    & 0.0   &  {\color{red}0.00097}\\
	d_{1\downarrow}^\dagger |&0.0   &  -0.24  & 0.0    & -0.04  & -0.00097 & 0.0\\
	d_{2\uparrow}^\dagger| &-0.04 &  0.0    & -0.24  & 0.0    & -0.04 &  0.0\\
	d_{2\downarrow}^\dagger| & 0.0    & -0.04  & 0.0     &-0.24  & 0.0   &  -0.04\\
	d_{3\uparrow}^\dagger| &0.0    & -0.00097  &-0.04 &  0.0 &    -0.24 &  0.0\\
	d_{3\downarrow}^\dagger| &0.00097 &  0.0   &  0.0 &    -0.04  & 0.0  &   -0.24	
\end{pmatrix}
\end{equation}

Betrachte z.B. das (1,6) Element ${\color{red}h_{16}}$, welches den Wert $h_{16} = 0.00097$ hat. Dies ist der Koeffizient von dem Term
\begin{equation}
	h_{16}d_{1\uparrow}^\dagger d_{3\downarrow}^\pdagger 
\end{equation} 
im elektronischen Teil des Hamiltonians, welcher verallgemeinert als 
\begin{equation}
	H_\mathrm{el} = \sum_{\vec{i}\vec{j}}h_{\vec{i}\vec{j}}^\pdagger d_\vec{i}^\dagger d_\vec{j}^\pdagger
\end{equation}
geschrieben werden kann mit $\vec{i}=(i,\sigma)$ ($i=\{1,2,3\}$ ist der Site-index und $\sigma=\{\uparrow,\downarrow\}$ der Spin-Index in diesem Beispiel).
Analog ist der Imaginärteil des elektronischen Hamiltonians zu lesen (btw, der Hamiltonian ist komplex, hoffe dass das kein Problem ist).
\subsection{*MODE$\alpha$*}
Diese Dateien sind die Koeffizienten der Terme der Elektron-Schwingungskopplung an die Mode $\alpha$. Beispiel (\texttt{hamiltonian\_1x3\_tPh\_-0.004\_thPd\_-0.004\_SOC0.0\_SOCPh0.001\_V0.52\_w00.002\_MODE1\_real.dat}) mit der Kopplungsmatrix an die erste Mode
\begin{equation}
	\mathrm{Re}(H_\mathrm{el-vib_1}) = \begin{pmatrix}
		& d_{1\uparrow} & d_{1\downarrow} & d_{2\uparrow} & d_{2\downarrow} & d_{3\uparrow} & d_{3\downarrow}\\ \hline
			d_{1\uparrow}^\dagger | & 0.0 &     0.0    & -0.004 &   0.0  &   0.0&    0.00097\\
			d_{1\downarrow}^\dagger | & 0.0   &  0.0    & 0.0    & {\color{red}-0.004} &   -0.00097 &  0.0\\
			d_{2\uparrow}^\dagger |& -0.004 &   0.0   &  0.0   &  0.0  &   0.0   &  0.0\\
			d_{2\downarrow}^\dagger |& 0.0    & -0.004&  0.0   &  0.0    & 0.0    & 0.0\\
			d_{3\uparrow}^\dagger |& 0.0    & -0.00097 &  0.0  &   0.0 &    0.0 &    0.0\\
			d_{3\downarrow}^\dagger |& 0.00097 & 0.0   &  0.0    & 0.0 &    0.0 &    0.0
	\end{pmatrix}(a_1^\dagger + a_1)
\end{equation}
Nehme als Beispiel das (2,4) Element ${\color{red} h_{14} = -0.004}$. Dies ist der Koeffizient des Terms 
\begin{equation}
	h_{14} d_{1\downarrow}^\dagger d_{2\downarrow} (a_1^\dagger + a_1)
\end{equation}
der Kopplung an die erste Schwingungsmode, also ganz allgemein sind in dieser Datei die Realteile des Koeffizienten der Terme 
\begin{equation}
	\sum_{\vec{i}\vec{j}} h_{\vec{i}\vec{j}}^\pdagger d_\vec{i}^\dagger d_\vec{j}^\pdagger(a_1^\dagger + a_1)
\end{equation}
Analog dazu sind in der Datei mit MODE2 im Dateinamen die Koeffizienten der Term
\begin{equation}
	\sum_{\vec{i}\vec{j}} h_{\vec{i}\vec{j}}^\pdagger d_\vec{i}^\dagger d_\vec{j}^\pdagger(a_2^\dagger + a_2).
\end{equation}
\end{document}




